% SPDX-FileCopyrightText: 2026 Schmiedmayer Lab and the project authors (see CONTRIBUTORS.md)
% SPDX-License-Identifier: MIT
\documentclass{article}
\begin{document}
\maketitle
\begin{abstract}
This illustrative manuscript exercises the Word review workflow with prose, citations, a table, a figure, and editable mathematics.
It does not report research findings.
\end{abstract}
\section{Methods}
\label{sec:methods}
The measured value was 10 units.

Notes for reviewers.
The template provides the workflow~\citep{schmiedmayer2026papertemplate}.
An inline equation $x^2 + y^2 = z^2$ remains editable in Word.
\begin{equation}
\label{eq:mean}
\bar{x} = \frac{1}{n}\sum_{i=1}^{n} x_i
\end{equation}
\section{Illustrative Outputs}
See \cref{sec:methods} for the procedure.
\begin{table}[ht]
\centering
\caption{Illustrative values for checking table formatting}
\label{tab:values}
\begin{tabular}{lrr}
\toprule
Condition & Value & Count \\
\midrule
Reference & 10.0 & 12 \\
Comparison & 12.5 & 12 \\
\bottomrule
\end{tabular}
\end{table}
\begin{figure}[ht]
\centering
\includegraphics[width=0.7\textwidth,alt={Four illustrative bars with increasing height}]{paper/figures/generated/review-bars.png}
\caption{Illustrative bars for checking image sizing and caption placement}
\label{fig:bars}
\end{figure}
\section{Declarations}
\subsection{Author Contributions}
\papercontributions
\bibliographystyle{unsrtnat}
\bibliography{references}
\end{document}
